%% kapa: exam=1 round=19 year=2008 period=1 form=A subject=accounting subject_ko=회계학 seq=9 number=89
%% source: Q-Net 감정평가사 기출문제 (cst00309) artlSeq=1001962 file=extracted/1008/1001962/2008년 제19회 감정평가사 1차시험 기출문제/1교시.hwp
%% conversion: kordoc 4.24.51 -> tools/gen1.py (auto); review=proofread-sonnet; flags=none
\begin{question}{89}
다음 중 무형자산의 회계처리에 대한 설명으로 가장 옳은 것은?
\begin{choices}
  \item 무형자산을 사용하는 동안 취득 후의 지출로 인하여 자산의 성능이 향상되는 경우에도 상각기간을 변경할 수 없다.
  \item 모든 무형자산은 상각기간의 제한이 없다.
  \item 무형자산의 상각은 당해 자산이 사용가능한 시점부터 상각한다.
  \item 무형자산은 정액법에 의해서만 상각한다.
  \item 무형자산의 상각이 다른 자산의 제조와 관련된 경우에는 판매비와관리비로 처리한다.
\end{choices}
\end{question}
\answer{3}
